\chapter{Derived Stack Construction for RSVP}

\section{Motivation and Overview}

In Chapter~5 the lamphrodynamic fields $(\Phi,\vec v,S)$ were introduced
as classical smooth fields on a Lorentzian spacetime $(M,g)$.
To quantize the theory using the AKSZ formalism,
we require a \emph{derived} moduli stack whose classical points correspond
to such triples, so that the mapping stack
\[
\mathrm{Map}(\mathsf{T}[1]M,\mathrm{Plenum})
\]
inherits a $(-1)$-shifted symplectic structure.
This chapter constructs the derived moduli stack $\mathrm{Plenum}$,
proves that it is a derived Artin stack locally of finite presentation,
describes its tangent complex and obstruction theory,
and prepares the ground for the shifted symplectic structure in
Chapter~7.

We rely on foundational results of Lurie \cite{Lurie-DAG} and
To\"en--Vezzosi \cite{ToenVezzosi-HAGII} on derived algebraic geometry,
and on Pantev--To\"en--Vaqui\'e--Vezzosi \cite{PTVV13} for the
construction of shifted symplectic forms on derived mapping stacks.


\section{The Derived Moduli Functor}

Fix a smooth oriented $4$-manifold $M$ with Lorentzian metric $g$.
For a derived affine scheme $\operatorname{Spec}A$,
define the space of lamphrodynamic fields over $A$ by
\[
\mathcal{U}(A)
:= \left\{
(\Phi,\vec v,S) \mid 
\Phi\in\Gamma(M_A,\mathcal{O}_{M_A}),
\;
\vec v\in\Gamma(M_A,T_{M_A}),
\;
S\in\Gamma(M_A,\mathcal{O}_{M_A})
\right\},
\]
where $M_A := M\times\operatorname{Spec}A$ and
$T_{M_A}$ is the relative tangent complex.
We regard $\mathcal{U}$ as a functor
\[
\mathcal{U} \colon \mathbf{dAff}^{\mathrm{op}}
\longrightarrow \mathbf{sSet}.
\]

\begin{definition}
The \emph{derived lamphrodynamic moduli stack} is the derived
stackification of $\mathcal{U}$, denoted
\[
\mathrm{Plenum} := \mathbf{L}\mathcal{U}.
\]
\end{definition}

Classical field configurations correspond to $A=\mathbb{R}$,
so that $\mathrm{Plenum}(\mathbb{R})$ recovers the classical triples
introduced in Chapter~5.


\section{Representability and Derived Artin Property}

We now establish that $\mathrm{Plenum}$ is a derived Artin stack.
Let us outline the representability argument; complete details follow
the general pattern of mapping stacks of sections of a vector bundle,
together with the derived enhancement required by the lamphrodynamic
constraint.

\begin{theorem}
\label{thm:Plenum-Artin}
The derived lamphrodynamic moduli stack $\mathrm{Plenum}$ is
a derived Artin stack locally of finite presentation over $\mathbb{R}$.
\end{theorem}

\begin{proof}[Proof (sketch)]
The moduli problem is locally given by sections $\Phi,S$ of trivial
line bundles and a section $\vec v$ of the tangent bundle.
Such mapping stacks are representable by derived Artin stacks
(\cite{ToenVezzosi-HAGII}, \S7).
The nonlinear constraints \eqref{eq:R-Phi}--\eqref{eq:compatibility}
are encoded as derived intersections of smooth Artin stacks,
preserving Artin representability. For details see
\cite{Lurie-DAG,ToenVezzosi-HAGII}.
\end{proof}

\begin{remark}
The lamphrodynamic constraint creates a derived structure:
classically one would impose $\mathcal{C}_\Phi=\mathcal{C}_v=0$,
but this would produce singularities in the moduli space.
Derived geometry resolves these singularities by enriching the
intersection with higher homotopical data.
\end{remark}


\section{Local Models and the Tangent Complex}

Given $\mathcal{U}$ as above,
the tangent complex at a point
\[
(\Phi_0,\vec v_0,S_0)\in \mathrm{Plenum}
\]
is given by the derived global sections
\begin{align}
T_{(\Phi_0,\vec v_0,S_0)}\mathrm{Plenum}
\simeq
\mathbf{R}\Gamma\bigl(
M,
\mathcal{O}_M \oplus T_M \oplus \mathcal{O}_M
\bigr).
\end{align}
When nonlinear constraints are included, the tangent complex becomes the
fiber of
\[
\mathbf{R}\Gamma(
M,
\mathcal{O}_M \oplus T_M \oplus \mathcal{O}_M
)
\longrightarrow
\mathbf{R}\Gamma(M, \mathcal{N}),
\]
where $\mathcal{N}$ is the derived normal bundle determined by the
Jacobian of the constraint functions
$\mathcal{C}_\Phi,\mathcal{C}_v$.

\begin{proposition}
The tangent complex of $\mathrm{Plenum}$ has perfect amplitude
contained in $[-1,1]$.
\end{proposition}

\begin{proof}[Proof (sketch)]
The tangent complex of sections of a perfect complex has perfect
amplitude given by the amplitude of the underlying complex of sheaves.
The derived constraint adds a fiber whose amplitude is bounded by $1$.
Standard arguments apply \cite{Lurie-DAG}.
\end{proof}

\begin{remark}
The amplitude $[-1,1]$ is crucial: it ensures $\mathrm{Plenum}$ has
a well-defined shifted symplectic structure of degree $2$; see Chapter~7.
\end{remark}

\section{Obstruction Theory and Derived Intersections}

We now analyze the derived enhancement induced by the nonlinear
lamphrodynamic constraints
\eqref{eq:R-Phi}--\eqref{eq:compatibility}.
Following \cite{ToenVezzosi-HAGII}, derived algebraic geometry replaces
classical intersections of moduli problems with homotopy fiber products.
Let $\mathcal{M}_{\Phi,S}$ denote the moduli stack of pairs of functions
$(\Phi,S)\in\Gamma(M,\mathcal{O}_M)^2$, and let $\mathcal{M}_{v}$ denote
the moduli stack of $\Gamma(M,T_M)$-valued sections.
The lamphrodynamic constraints define classical morphisms
\[
\mathcal{M}_{\Phi,S}\times \mathcal{M}_{v}
\longrightarrow \mathcal{N}_{\Phi,v}
\]
whose zero locus would be singular in general.

\begin{definition}
The \emph{derived lamphrodynamic intersection} is the homotopy fiber
product
\[
\mathrm{Plenum}
\simeq
(\mathcal{M}_{\Phi,S}\times \mathcal{M}_{v})
\times^{\mathbf{R}}_{\mathcal{N}_{\Phi,v}} \ast.
\]
\end{definition}

\begin{remark}
The derived homotopy fiber product retains information lost by
classical intersection, encoding obstructions in higher homotopical
degrees. This derived structure is responsible for the shifted
symplectic geometry we exploit in Chapter~7.
\end{remark}


\section{Derived Normal Bundle and Obstruction Classes}

Let $\mathcal{N}_{\Phi,v}$ denote the normal bundle determined by the
Jacobian of the lamphrodynamic constraint functions.
The derived normal bundle is obtained by taking the homotopy fiber
of the Jacobian mapping between perfect complexes, producing a derived
vector bundle in degrees $[-1,0]$.
The obstruction theory for $\mathrm{Plenum}$ is hence controlled by the
derived cohomology of $\mathcal{N}_{\Phi,v}$, localized on $M$.

\begin{theorem}
\label{thm:obstruction-Plenum}
The derived moduli stack $\mathrm{Plenum}$ admits a perfect obstruction
theory in the sense of \cite{BF97}, induced by the derived normal bundle
of the lamphrodynamic constraint.
\end{theorem}

\begin{proof}[Proof (sketch)]
The lamphrodynamic constraint defines a morphism of derived Artin stacks.
The homotopy fiber of the Jacobian defines a perfect obstruction theory,
using the general representability results of
\cite{BF97,ToenVezzosi-HAGII}.
\end{proof}

\begin{remark}
Physically, the obstruction theory describes infinitesimal lamphrodynamic
defects: perturbations of $(\Phi,\vec v,S)$ that are obstructed by the
constraint. In Volume~II these play a role near entropic singularities.
\end{remark}


\section{Existence of a $2$-Shifted Symplectic Structure}

We now show that $\mathrm{Plenum}$ admits a natural $2$-shifted
symplectic structure. The general existence theorem in \cite{PTVV13}
applies to derived intersections of moduli of sections of perfect
complexes when certain orientation data are available. In the present
setting, the Lorentzian manifold $(M,g)$ provides the necessary
orientation and volume form, and the kinetic and coupling terms in the
action provide a natural closed $3$-form.

Let $\Theta$ be the lamphrodynamic $3$-form defined locally by
\begin{equation}
\label{eq:local3form}
\Theta :=
\Phi\, dS\wedge \star d\Phi
+ \beta\, \langle\nabla\Phi,\nabla S\rangle\, dV_g
+ \gamma\, \langle\nabla S,\vec v\rangle\, d(\iota_{\vec v}dV_g),
\end{equation}
where $\star$ is the Hodge operator associated to $g$ and
$\gamma$ is a normalizing constant determined below.

\begin{proposition}
The form $\Theta$ extends to a closed $3$-form on the derived moduli
stack $\mathrm{Plenum}$, defining a $2$-shifted presymplectic structure.
\end{proposition}

\begin{proof}[Proof (sketch)]
Each term in \eqref{eq:local3form} is closed at the classical level
modulo the lamphrodynamic constraint. The derived extension is obtained
by pulling back to the derived intersection and verifying closedness in
the derived de~Rham complex, following \cite{PTVV13}. Full details will
be provided in Appendix~E.
\end{proof}


\section{Local Charts and Explicit Expression for $\omega$}

We now identify the $2$-shifted symplectic form $\omega$ as the derived
differential of $\Theta$:
\[
\omega := d_{\mathrm{dR}}\Theta.
\]
Classically, this yields
\begin{align}
\omega
&=
d\Phi\wedge dS\wedge \star d\Phi
+ \Phi\, d(dS\wedge \star d\Phi)
\\
& \qquad
+ \beta\, d\bigl(
\langle\nabla\Phi,\nabla S\rangle\, dV_g\bigr)
+ \gamma\, d\bigl(
\langle\nabla S,\vec v\rangle\, d(\iota_{\vec v}dV_g)
\bigr).
\end{align}
The derived extension incorporates infinitesimal deformations and ghosts
of each lamphrodynamic field. We omit the full expression here, as it
requires the total derived de~Rham complex on an Artin stack; see
Appendix~E.

\begin{theorem}[Shifted symplectic structure of $\mathrm{Plenum}$]
\label{thm:shifted}
The closed derived $3$-form $\Theta$ induces a
$2$-shifted symplectic structure $\omega$ on $\mathrm{Plenum}$.
\end{theorem}

\begin{proof}[Proof (sketch)]
Closedness is by construction. Non-degeneracy follows from the perfect
amplitude $[-1,1]$ of the tangent complex and the existence of
the lamphrodynamic constraint which produces a derived intersection of
section spaces satisfying the orientation axioms of \cite{PTVV13}.
See also \cite{ToenVezzosi-HAGII}.
\end{proof}

\begin{remark}[Physical origin of the shift]
Since the physical spacetime dimension is $4$,
integration of the derived $3$-form over the source $M$ produces an
effective shift by $3-4=-1$ in the BV symplectic degree, giving a
$(-1)$-shifted BV structure in Chapter~7. Hence the physical degree $k$
for RSVP quantization is forced to be $2$ at the classical level,
consistent with the AKSZ construction for $4$-dimensional theories.
\end{remark}

\section{Main Structural Theorem for the Lamphrodynamic Plenum}

We now unify the constructions of the preceding sections into a single
structural result, collecting assumptions and conclusions in one place.

\begin{theorem}[Main Structural Theorem]
\label{thm:Main-Plenum}
Let $M$ be a smooth oriented time-oriented $4$-dimensional Lorentzian
manifold and let $(\Phi,\vec v,S)$ be the lamphrodynamic fields with
nonlinear constraints \eqref{eq:R-Phi}--\eqref{eq:compatibility}.
Then the following properties hold:
\begin{enumerate}
\item The moduli problem of lamphrodynamic triples extends to a derived
functor $\mathcal{U}\colon\mathbf{dAff}^{\mathrm{op}}\to\mathbf{sSet}$
whose derived stackification $\mathrm{Plenum}$ is a derived Artin stack
locally of finite presentation.

\item The nonlinear constraints are resolved by a perfect obstruction
theory induced from the derived normal bundle of the constraint, and
$\mathrm{Plenum}$ has a tangent complex of perfect amplitude $[-1,1]$.

\item There exists a canonical closed derived $3$-form $\Theta$ whose
derived de~Rham differential $\omega := d_{\mathrm{dR}}\Theta$
defines a $2$-shifted symplectic structure on $\mathrm{Plenum}$.

\item The shifted symplectic structure is functorial with respect to
pullback along smooth morphisms of Lorentzian spacetimes,
and admits local presentations on Artin atlases of $\mathrm{Plenum}$
compatible with the lamphrodynamic constraint.
\end{enumerate}
\end{theorem}

\begin{proof}[Proof (sketch)]
Part (1) follows from representability of mapping stacks of sections of
perfect complexes and stability of Artin representability under derived
homotopy fiber products \cite{ToenVezzosi-HAGII}.
Part (2) follows from the construction of the derived normal bundle and
Behrend--Fantechi obstruction theories \cite{BF97}.
Part (3) is an instance of the PTVV construction of shifted symplectic
forms on derived mapping stacks with perfect obstruction theory
\cite{PTVV13}.
Part (4) follows from standard functoriality properties of shifted
symplectic geometry \cite{Calaque16}.
\end{proof}


\section{Physical Interpretation}

Although the shift by $2$ appears purely technical, it reflects the fact
that RSVP is a $4$-dimensional physical theory whose classical degrees
of freedom arise from a 3-form $\Theta$. The descent from $k=2$ to $k-4=-2$
under integration over the source would ordinarily yield a $(-2)$-shift,
but the AKSZ construction uses $\mathsf{T}[1]M$ in place of $M$,
producing the canonical $(-1)$-shift required for the BV bracket.
Thus the appearance of degree $2$ is a physical consequence of the
lamphrodynamic structure in four dimensions.

\begin{remark}
One may view the lamphrodynamic plenum as a universal ambient space of
physical fields in which gravitational and entropic interactions are
encoded as derived symplectic geometry. In this sense, $(\Phi,\vec v,S)$
constitute a pre-geometric substrate from which spacetime causality
emerges (developed further in Volume~II).
\end{remark}


\section{Preview of the AKSZ--RSVP Action}

In Chapter~7 we construct the AKSZ action associated to the
$2$-shifted symplectic derived stack $\mathrm{Plenum}$.
The result will be a $(-1)$-shifted symplectic BV theory on the
mapping stack $\mathrm{Map}(\mathsf{T}[1]M,\mathrm{Plenum})$,
providing a geometric and functorial route to quantization.
The Euler--Lagrange equations of the AKSZ action reproduce the
lamphrodynamic field equations of Chapter~5, up to additional ghost and
antifield terms required for BV consistency.

\begin{remark}
In the AKSZ formalism, no gauge-fixing or ad-hoc ghosts are introduced:
the ghost and antifield sectors appear automatically, determined by the
shifted symplectic geometry. This is essential for a physically
well-defined quantization of RSVP.
\end{remark}
